% l-analisis_de_resultados.tex
% !TEX root = z-main.tex

\section{Análisis de resultados}

Este capítulo discute el conjunto de modelos desarrollados (Montículos 3, 7, 12, 13, 14 y La Palangana), analizando el flujo de trabajo empleado y los resultados alcanzados en tres etapas: modelado geométrico, aplicación de materiales/texturas e integración prevista a un ambiente de realidad aumentada (RA). Se incluyen, además, consideraciones de validación arqueológica, rendimiento, limitaciones y líneas de trabajo futuro.

\subsection{Modelado geométrico}

Se completó el modelado tridimensional de los seis montículos propuestos siguiendo un procedimiento uniforme. El flujo por etapas (malla $\rightarrow$ sólido $\rightarrow$ texturizado) permitió validar progresivamente la calidad del resultado:
\begin{itemize}
    \item \textbf{Geometría base (wireframe).} Esta vista facilitó la inspección de proporciones, alineación de escalones, continuidad de taludes y simetrías. Fue útil para detectar y corregir aristas/caras innecesarias, evitar geometría y asegurar una topología limpia.
    \item \textbf{Volumetría (sólido).} La visualización sin materiales aisló la forma general de cada montículo y ayudó a comprobar la relación entre plataformas, escalinatas y superestructuras, así como la pendiente y los cambios de nivel.
    \item \textbf{Aplicación de materiales (texturizado).} En esta fase se verificó la correspondencia entre las texturas asignadas y los materiales originales documentados, cuidando la orientación y escala para obtener una representación realista.

\end{itemize}

\subsection{Texturizado y materiales}

La aplicación de materiales se orientó a reflejar acabados de barro, paja y maderas creíbles para Kaminaljuyú, manteniendo uniformidad visual entre modelos:
\begin{itemize}
    \item \textbf{Criterios visuales.} Se priorizaron gamas de barro en superficies expuestas, y tonos rocosos para bordes o elementos constructivos, siguiendo recomendaciones del área de Arqueología.
    \item \textbf{UV y repetición.} Se cuidó el desplegado UV para minimizar estiramientos perceptibles y favorecer patrones repetibles sin artefactos notorios en planos amplios (plataformas y taludes).
    \item \textbf{Preparación para RA.} La cantidad de materiales se mantuvo acotada para reducir mallas, y las texturas pueden consolidarse y emplear compresión en la exportación.
\end{itemize}

\subsection{Modelos en ambiente virtual (RA)}

Los modelos tridimensionales fueron preparados para su integración en un entorno de realidad aumentada en Android, en el cual podrán visualizarse superpuestos al espacio físico. Para facilitar su uso, cada modelo se exportó en formato \emph{.obj} con su correspondiente archivo \emph{.mtl}, manteniendo materiales simples y compatibles con los motores gráficos empleados en la aplicación. Se cuidó la orientación en los ejes y la ubicación del origen en (0,0,0), de modo que los modelos puedan escalarse y posicionarse con facilidad dentro del entorno virtual. Asimismo, se verificó la coherencia de escala entre los seis montículos y la uniformidad del estilo visual, garantizando una correcta correspondencia espacial y visual al momento de su despliegue en RA. Aunque la implementación de la aplicación corresponde a otro integrante del equipo, los modelos se entregaron optimizados y listos para su uso en dicha plataforma.

\subsection{Validación arqueológica y coherencia entre modelos}

El conjunto mantiene una lectura coherente con los rasgos registrados para Kaminaljuyú:
\begin{itemize}
    \item \textbf{Forma escalonada.} Las proporciones de plataformas y la relación con la escalinata principal son consistentes entre montículos, respetando las variaciones particulares observadas en cada caso.
\end{itemize}

\subsection{Comparación y síntesis del conjunto}

En conjunto, los seis modelos mantienen una coherencia formal y visual que refleja adecuadamente las características arquitectónicas de Kaminaljuyú. Todos fueron elaborados mediante el mismo flujo de trabajo, lo que permitió conservar la proporción entre plataformas, taludes y escalinatas, así como una topología limpia y eficiente gracias a la eliminación de caras no visibles. El resultado evidencia que el pipeline es reproducible y puede aplicarse a futuros montículos o variantes sin comprometer la consistencia general.

En términos de balance entre fidelidad y rendimiento, la simplificación geométrica se mantuvo dentro de límites adecuados para conservar los rasgos constructivos esenciales, garantizando al mismo tiempo un desempeño óptimo para su uso en entornos de realidad aumentada. Asimismo, la uniformidad de materiales y texturas, con excepción del Montículo 14, que requirió ajustes particulares, favorece la lectura visual del conjunto como parte de un mismo sitio arqueológico.

Finalmente, la representación tridimensional obtenida contribuye a propósitos educativos y de divulgación, al ofrecer una forma accesible y visualmente coherente de aproximarse a cómo pudieron verse originalmente las estructuras, facilitando su comprensión por parte del público general y de visitantes en experiencias de realidad aumentada.

\subsection{Limitaciones}

\begin{itemize}
    \item \textbf{Datos de partida.} Algunas suposiciones de materiales y acabado superficial son creíbles pero no definitivas, podrían refinarse con campañas de documentación adicionales.
    \item \textbf{Iluminación de referencia.} Las vistas de \emph{viewport} no equivalen a un render físicamente basado, el aspecto puede variar con condiciones de iluminación realistas en RA.
    \item \textbf{Detalle fino.} Elementos menores (erosión, irregularidades, grietas) fueron simplificados para priorizar legibilidad y rendimiento.
\end{itemize}

\subsection{Trabajo futuro}

\begin{itemize}
    \item \textbf{Afinar materiales PBR.} Incorporar mapas de normales y oclusión ambiental derivados de texturas reales o fotogrametría puntual para detalles finos.
    \item \textbf{Ampliación del corpus.} Extender el pipeline a otros montículos o variantes históricas y documentar discrepancias para análisis comparativo.
\end{itemize}

\section{Discusión de resultados}

Ver si se discute solamente lo de la encuesta o también lo de los resultados del modelado.